\section{Technical Moat}
\label{sec:moat}

\subsection{The compounding asset stack}

\begin{diagram}
   hardware knowledge            (MCU/PHY/transceiver facts, per-part
        +                         init sequences, errata workarounds)
   vehicle topology knowledge    (per-OEM-platform network layouts,
        +                         gateway routing patterns)
   semantic peripheral ontology  (the closed set of device classes and
        +                         their capability contracts, Section 2)
   verified drivers              (Verified Primitive Library entries,
        +                         Section 5)
   verified application          (tested Behavior-layer compositions
   primitives                     for common functions: door-lock
        +                         interlocks, lighting state machines)
   network mappings              (Vehicle Knowledge Graph entries,
        +                         confidence-scored, reusable within
                                   and often across same-platform
                                   vehicles)
   firmware generation           (Target Packs, resource-planner
        +                         constraint models per MCU family)
   hardware measurements         (real electrical characterization
        +                         data: actual current draw, actual
                                   servo bandwidth, actual EMI behavior
                                   of specific supplier peripherals)
   validation evidence           (SIL/HIL results, static analysis
                                   history, certification-evidence
                                   bundles per primitive/target)
\end{diagram}

Every layer above is populated by \emph{doing the work} of bringing one more machine under Alfred, and every layer above is reusable on the next machine to a degree a point-solution competitor's equivalent artifact is not: a CAN analyzer vendor's decoded-signal database is not attached to a verified, deployable driver; a firmware-generation tool's templates are not attached to real measured electrical behavior of the specific parts a customer actually uses; an RTOS vendor's board-support packages are not attached to a semantic ontology that lets an application call \code{pump.set\_speed()} without knowing which BSP is underneath.

\subsection{What compounds fastest, and why}

\begin{itemize}
\item \textbf{Hardware Knowledge Graph entries for common parts} (a specific CAN transceiver's init quirks, a specific motor driver IC's fault-reporting behavior) are reusable across every customer and every vehicle class immediately --- this is the highest-leverage layer to invest in early because its reuse rate is nearly 100\%.
\item \textbf{Per-OEM-platform network mapping patterns} (Option B) compound within a platform family fast (a second vehicle from the same OEM/platform generation shares most CAN-ID semantics with the first) and compound more slowly but still meaningfully across platforms (turn-signal-style signals tend to occupy structurally similar positions across many OEMs' body networks for historical/supplier reasons) --- this is genuinely proprietary data no competitor can buy, because it is built from Alfred's own accumulated experiment/correlation runs, not from public specs.
\item \textbf{Verified Primitive Library entries} compound across vehicle classes, not just within one: a verified brushed-DC-motor-with-Hall-feedback primitive built for an automotive seat motor is directly reusable for an agricultural implement actuator or a UGV control surface, because the semantic ontology (Section~2) was deliberately designed around electrical interface types, not vehicle-specific function names.
\item \textbf{Validation/certification evidence} compounds specifically for regulated programs: a primitive that has already accumulated a track record of HIL results and static analysis history across N prior deployments is cheaper to qualify for the N+1th program's functional-safety case than a from-scratch primitive, which becomes a genuine switching cost for a customer once they have several primitives in production.
\end{itemize}

\subsection{5--10 year moat candidates}

\begin{enumerate}
\item \textbf{The cross-vehicle Hardware Knowledge Graph itself}, once it spans enough parts and platforms that a competitor starting today would need years of equivalent field engagements to reach parity --- this is a data moat in the classic sense: value scales with usage, and usage is gated by having already done the (expensive, physical, non-scrapable) work of touching real machines.
\item \textbf{The Verified Primitive Library's breadth and track record}, which functions like a certified-parts catalog that gets more valuable to a safety-conscious customer the longer its field history is, similar in kind to how an RTOS's certification pedigree (e.g., a pre-certified SafeRTOS kernel) becomes a switching cost independent of the code itself.
\item \textbf{The deterministic firmware-generation pipeline's Target Pack ecosystem}, if Alfred reaches enough MCU-family coverage that "which MCU can Alfred already generate firmware for" starts to shape customers' own hardware choices, analogous to how AUTOSAR MCAL vendor support shapes OEM MCU selection today.
\item \textbf{Structured-experiment-derived semantic mapping patterns for closed/proprietary industrial and ag/mining control systems}, where no public DBC/ARXML equivalent exists at all --- this is the category where Alfred's discovery methodology is hardest to replicate quickly because there is no shortcut through public documentation; it must be earned per platform through the correlation/experiment pipeline.
\end{enumerate}

None of these moats exist at MVP0--2 scale (Section~\ref{sec:roadmap}); they are the reason the roadmap prioritizes getting real engagements running as early as possible even with a minimal feature set --- the moat is a function of elapsed engagements and accumulated verified artifacts, not of how complete the initial architecture is.
